% Generated draft: requires mathematical and source review.
\documentclass[11pt]{article}
\usepackage[margin=25mm]{geometry}
\usepackage{fontspec}
\setmainfont{Latin Modern Roman}
\usepackage{amsmath,amssymb,mathtools,unicode-math}
\setmathfont{Latin Modern Math}
\usepackage{enumitem,needspace,xurl,hyperref}
\setlength{\parindent}{0pt}
\setlength{\parskip}{5pt}
\setlength{\emergencystretch}{3em}
\newcommand{\N}{\mathbb{N}}
\newcommand{\Z}{\mathbb{Z}}
\newcommand{\Q}{\mathbb{Q}}
\newcommand{\R}{\mathbb{R}}
\newcommand{\C}{\mathbb{C}}
\newcommand{\F}{\mathbb{F}}
\newcommand{\A}{\mathbb{A}}
\newcommand{\PP}{\mathbb{P}}
\newcommand{\E}{\mathbb{E}}
\newcommand{\eps}{\varepsilon}
\newcommand{\dd}{\,\mathrm{d}}
\begin{document}
\section*{UnsolvedMath ID: 800627}
\textbf{Source problem code:} catalogue entry
\subsection{Definitions and mathematical statement}
Let \(\mathbb N=\{1,2,\ldots\}\). For an abelian topological group \(G\), let \(\mathcal M(G)\) denote its meager subsets, where a set is meager if it is a countable union of nowhere dense sets. When a measure on \(G\) is specified, let \(\mathcal N(G)\) denote the sets of outer measure zero. For \(\mathcal I\in\{\mathcal M,\mathcal N\}\), call \(A\subseteq G\) \(\mathcal I\)-additive if
\[
\forall B\subseteq G\quad
\bigl(B\in\mathcal I(G)\Longrightarrow A+B\in\mathcal I(G)\bigr),
\qquad A+B=\{a+b:a\in A,\ b\in B\}.
\]
For \(G=\mathbb R\), use ordinary addition, the usual topology, and Lebesgue measure. The standard interpretation for \(C=\{0,1\}^{\mathbb N}\) uses the product topology, coordinatewise addition modulo \(2\), denoted by \(\oplus\), and the product probability measure assigning probability \(1/2\) to each coordinate value. These ambient-group and measure conventions make the usual meanings explicit; the supplied record does not itself define them.

The source asks whether meager-additivity and, separately, null-additivity transfer from \(\mathbb R\) to \(C\) for a set \(X\subseteq[0,1]\). Its phrase ``when viewed as a subset of \(\{0,1\}^{\mathbb N}\)'' leaves the identification unspecified. To display the intended implications precisely conditional on a binary-representation choice, put
\[
\pi:C\longrightarrow[0,1],\qquad
\pi(s)=\sum_{n=1}^{\infty}s_n2^{-n},
\]
and let \(b:[0,1]\longrightarrow C\) be a fixed choice of binary expansion satisfying \(\pi(b(x))=x\) for every \(x\in[0,1]\). For this interpretation, the two questions are whether, for every \(X\subseteq[0,1]\), respectively,
\[
\begin{aligned}
&\left[\forall M\in\mathcal M(\mathbb R),\ X+M\in\mathcal M(\mathbb R)\right]
\Longrightarrow
\left[\forall K\in\mathcal M(C),\ b(X)\oplus K\in\mathcal M(C)\right],\\
&\left[\forall N\in\mathcal N(\mathbb R),\ X+N\in\mathcal N(\mathbb R)\right]
\Longrightarrow
\left[\forall L\in\mathcal N(C),\ b(X)\oplus L\in\mathcal N(C)\right].
\end{aligned}
\]
Here \(b(X)=\{b(x):x\in X\}\). The two implications are separate questions, not a joint hypothesis requiring both kinds of additivity. The record does not specify a binary-expansion convention, whether the intended realization is instead the full inverse image \(\pi^{-1}(X)\), or a quantifier over representation choices. Thus the displayed formulation is conditional on the stated interpretation and requires source clarification before being identified with the exact intended formulation.
\subsection{Short English statement}
For a set \(X\subseteq[0,1]\), does the property that adding any meager real set gives a meager set persist when \(X\) is represented in the binary sequence group? Ask the same question separately for null sets. The representation and ambient-group conventions require confirmation. The registry labels this problem open; that status is supplied metadata, not an independently verified conclusion.
\subsection{Sources}
\begin{itemize}
\item Supplied local UnsolvedMath registry record: permanent ID \(800627\), exact source problem code \texttt{catalogue entry}, title ``Selection Principles in Mathematics --- Problem 56 ([prods]).'' The record lists Tsaban as proposer, gives the year 2006, and labels the problem \texttt{open}.
\item Supplied source attribution: Tsaban, \emph{Selection Principles and Special Sets of Reals: Open Problems} (2006), Problem 56. Supplied URL: \url{https://arxiv.org/abs/math/0609601}. The background records access on 2026-10-16 and extraction by \texttt{arxiv\_latex}; these are supplied provenance statements, not access or verification performed for this draft.
\item The supplied attribution/context marker is \texttt{[prods]}. Its bibliographic expansion is absent from the record and has not been reconstructed. The supplied upstream revision is \texttt{d371ee917ac9647fbb06cd6fc529deec6f3dda82}.
\end{itemize}
\end{document}
